\chapter{Random numbers library}
\label{chap:ran}

\section{Introduction}

\subsection{Supported distributions}

Distributions that are currently supported are summarized in
Table~\ref{tab:distributions}.
At each distribution is associated a numerical constant. The total number
of distribution can be obtained by taking \verb|RAN_CONSTANT_END|-1,
and a constant \verb|RAN_CONSTANT| is also defined.
A multivariate distribution can be identified with the prefix
\verb|m_| before the distribution name.

The number of parameters associated to some distribution can be
obtained by postfix \verb|_NPAR| to the numerical constant name. The
library supports various continuous distributions, as well as some
discrete ones.
This is however only valid for univariate distributions since the
number of parameters for multivariate distributions depends on the
vector dimension.

\begin{table}[htbp]
\begin{center}
\begin{tabular}{|ll|l|}
\hline
{\bf Distribution} & {\bf Identifier} & {\bf \#Parameters} \\
\hline
\multicolumn{3}{|l|}{\bf Univariate distributions} \\
\multicolumn{3}{|c|}{\textit{Discrete distributions}} \\
Bernouilli & \verb|RAN_BERNOUILLI| & 1 \\
Dirac & \verb|RAN_DIRAC| & 1 \\
Geometric & \verb|RAN_GEOMETRIC| & 1 \\
Integer in $[0, \mbox{max}[$ & \verb|RAN_INTEGER| & 1 \\
\multicolumn{3}{|c|}{\textit{Continuous distributions}} \\
Arcsinus & \verb|RAN_ARCSINE| & 0 \\
Cauchy & \verb|RAN_CAUCHY| & 2 \\
%Erlang \\
Exponential & \verb|RAN_EXPONENTIAL| & 1 \\
Gamma & \verb|RAN_GAMMA| & 2 \\
Gumbel & \verb|RAN_GUMBEL| & 2 \\
Johnson & \verb|RAN_JOHNSON| & 4 \\
Laplace & \verb|RAN_LAPLACE| & 2 \\
Lognormal & \verb|RAN_LOGNORMAL| & 2 \\
Normal & \verb|RAN_NORMAL| & 2 \\
Student & \verb|RAN_STUDENT| & 1 \\
Triangular & \verb|RAN_TRIANGULAR| & 0 \\
Truncated normal & \verb|RAN_TRUNCATED_NORMAL| & 4 \\
Truncated normal B & \verb|RAN_TRUNCATED_NORMAL_B| & 4 \\
Uniform on (0,1) & \verb|RAN_UNIFORM| & 0 \\
\hline
\multicolumn{3}{|l|}{\bf Multivariate distributions} \\
\multicolumn{3}{|c|}{\textit{Continuous distributions}} \\
Multivariate normal & \verb|RAN_M_NORMAL| \\
Unit sphere & \verb|RAN_M_UNIT_SPHERE| \\
\hline
\end{tabular}
\end{center}
\label{tab:distributions}
\caption{Supported distributions.}
\end{table}

\subsection{Terms and symbols}

We use these notational conventions:
\begin{itemize}
\item
$\mathcal{N}$: natural set.
\item
$\mathcal{N}_0$: natural set, with the exclusion of 0.
\end{itemize}

%%%%%%%%%%%%%% UNIFORM DISTRIBUTION %%%%%%%%%%%%%%

\section{Uniform random numbers generation}
\label{sec:uniform}

A good uniform random generator on the interval $[0,1]$ is the major
component of any good random generator library. Draws from other
distributions are usually obtained by adequately transform an
uniformly distributed sample (see Levroye~\cite{Levr86} for a survey
of these techniques).
Many researches have been conducted to develop good random generators,
and lead to different algorithms.
The search is however far from be over and
improved generators will probably be released (see for instance
L'Ecuyer~\cite{LEcu97}).
The reader is therefore invited to include such algorithms if she/he feels the need to do it.
Algorithms can be tested with the package TestU01 (L'Ecuyer et
Simard~\cite{LEcuSima07}).

The basic principle of a uniform random numbers generator is to define
a transition function $f: \mathcal{S} \rightarrow \mathcal{S}$, where
$\mathcal{S}$ is the state space.
The cardinality of $\mathcal{S}$ is assumed to be finite.
The initial state is denoted by $s_0$, and we will write
\[
s_n = f(s_{n-1}).
\]
We will furthermore assume that $f$ is periodic for all $n$ greater of
equal to some known $\tau$ (often equal to 0), with the period denoted
by $\rho$.
In other terms, we have $s_{n+\rho} = s_n$ for all $n\ge\tau$.
We denote the output space by $\mathcal{U}$, and we assume here that
$\mathcal{U} = (0,1)$.
The output function $g: \mathcal{S} \rightarrow \mathcal{U}$
transforms the state $s_n$ into an output value $u_n$.
This behavior is graphically illustrated in Figure~\ref{fig:rng}, with
$\tau = 0$.
\begin{figure}[htbp]
\[
\begin{CD}
 \cdots @>f>> s_{\rho-1} @>f>> s_0 @>f>> s_1  @>f>> 
  \cdots @>f>> s_n @>f>> \cdots \\ 
%
  @. @V{g}VV  @V{g}VV   @V{g}VV   
  @.  @V{g}VV  \\
%
 \cdots @. u_{\rho-1} @. u_0
 @.  u_1  @.   \cdots @.  u_n @.  \cdots \\
\end{CD}
\]
\label{fig:rng}
\caption{Behavior of a uniform random numbers generator}
\end{figure}

The basic function is
\begin{verbatim}
Random *ran_random_new(int id, unsigned long *seed, int n);
\end{verbatim}
This create an object of type \verb|Random *|.
\verb|seed| is a vector of integers that will be used to initialize
the state of the random numbers generator, that is dependent of the
underlying $U(0,1)$ generator (see Section~\ref{sec:uniform}).
If the seed is set to NULL, a default seed will be used.
The length of this vector also depends of the $U(0,1)$-generator, but
each generator can be initialized with a seed dimension equal to 1.
If no specific random numbers generator is specified, i.e. if
\verb|id| is set to 0, the
\verb|RAN_DEFAULT| generator is used.
A specific random numbers generator can be obtained by using the
identifier given in Table~\ref{tab:generators}.
These generators are discussed in more details in the rest of the
document.
\begin{table}[htb]
\begin{center}
\begin{tabular}{|c|c|c|}
\hline
{\bf Name} & {\bf Identifier} & {\bf Seed size} \\
\hline
Default generator & \verb|RAN_DEFAULT| & \verb|RAN_DEFAULT_SEED_LENGTH| \\
\hline
Standard minimal & \verb|RAN_STANDARDMINIMAL| & 1 \\
\hline
MRG32K3a & \verb|RAN_MRG32K3A| & 6 \\
\hline
MT19937 & \verb|RAN_MT19937| & $[1,\ldots,624]$\\
\hline
\end{tabular}
\caption{uniform random numbers generators}
\label{tab:generators}
\end{center}
\end{table}

A simplified random numbers constructor is
\begin{verbatim}
Random *ran_random();
\end{verbatim}
which is equivalent to \verb|ran_random_new(0, NULL, 0);\verb|
A copy of the initial seed is stored inside the generator, which can
be retrieved with the function
\begin{verbatim}
unsigned long *ran_random_get_seed(const Random *random);
\end{verbatim}
You can also change this seed using the following function:
\begin{verbatim}
void ran_random_set_seed(Random *random, unsigned long *seed, int n)
\end{verbatim}
The rng internal seed is stored independently of the given parameter seed.

The state structure is specific to each uniform random generator, so
the generic function uses the type \verb|void *|.
The reader should refer to the corresponding generator to learn which
structure has to used.
If this pointer is \verb|NULL|, an arbitrary initial state will be
generated using the internal clock, and this value is place in the
address pointed by \verb|state|.
The state can also be retrieved with the function
\begin{verbatim}
void *ran_random_get_state(Random *random);
\end{verbatim}

The memory associated to the generator has to be freed with the
function
\begin{verbatim}
void ran_random_free(Random *random);
\end{verbatim}

The name of the currently used generator can be obtained with the
function
\begin{verbatim}
char *ran_random_generator_name(Random *random);
\end{verbatim}
You can use a different algorithm for
random generation, by calling the function
\begin{verbatim}
void random_set_generator(Random *random, int id,
                          unsigned long *seed, int n)
\end{verbatim}
\verb|id| is a constant describing the generator to use and seed is an
integers vector of dimension $n$, used to define the initial state, as
described in Section~\ref{sec:uniform}.
Table~\ref{tab:generators} gives supported generators and the
corresponding constants, as well as the corresponding seed vector size.

\subsection{Linear Congruential Generators}
\index{generator!linear congruential}

Linear congruential generators (LCGs) have been introduced by
Lehmer~\cite{Lehm51}.
A sequence of integers $Z_1, Z_2, \ldots$ is defined by the recursive
formula
\begin{equation}
Z_i = (aZ_{i-1}+c) \mod m.
\label{eq:lcg}
\end{equation}
A first desirable property of a LCG is that it has full period.
The following theorem, proved by Hull and Dobell~\cite{HullDobe62}, gives
such a characterization.

\begin{thm}
The LCG defined in \eqref{eq:lcg} has full period if and only if the
following three conditions hold:
\begin{enumerate}
\item
the only positive integer that (exactly) divides both $m$ and $c$ is
1;
\item
if $q$ is a prime number that divides $m$, then $q$ divides $a-1$;
\item
if 4 divides $m$, then 4 divides a-1.
\end{enumerate}
\end{thm}

If $c = 0$, a period equal to $m$ cannot be achieved, but under some
conditions, we can obtain a period of $m-1$.
For $m$ prime, it can be shown that the period is $m-1$ if $a$ is a
primitive element modulo m, that is the smallest integer $l$ for which
$a^l-1$ is divisible by $m$ is $l=m-1$ (Knuth, 1988, p20).
LCGs that fullfil these conditions are called prime modulus
multiplicative LCGs (PMMLCGs).

A generic linear congruential generator can be created using the function
\begin{verbatim}
LCG *ran_lcg_new(unsigned long seed);
\end{verbatim}
where seed is the initial state of the generator.
The allocated memory memory can be released using the fonction
\begin{verbatim}
void ran_lcg_free(void *lcg);
\end{verbatim}

\begin{verbatim}
LCG *ran_lcg_new_with_state(unsigned long *state);
LCG *ran_lcg_new_with_seed(unsigned long seed);
unsigned long *ran_lcg_get_state(void *lcg);
void ran_lcg_set_state(void *lcg, void *state);
void ran_lcg_init(Random *random, unsigned long seed);
\end{verbatim}

By default, the created generator is the Standard Minimal generator (see
Section~\ref{sec:standard_minimal}, but it the parameters can be
modified using the function
\begin{verbatim}
void ran_lcg_set_parameters(LCG *lcg, double a, double c, double m);
\end{verbatim}
and a draw can be obtained with the function
\begin{verbatim}
double ran_lcg_get_val(void *lcg);
\end{verbatim}
The following code defines the infamous \verb|RANDU| generator:
\begin{lstlisting}
  unsigned long m = 2;

  r = ran_lcg_new(seed);
  for (i = 0; i < 30; i++) m *= 2;
  ran_lcg_set_parameters(r, 65539, 0, m);
\end{lstlisting}

\subsubsection{Standard minimal generator}
\index{generator!standard minimal}
\label{sec:standard_minimal}

A popular, but outdated random generator is the standard minimal,
following the terminology introduced by Park and
Miller~\cite{ParkMill88}, while it is known to be not optimal
(L'Ecuyer and Simard~\cite{LEcuSima07}).
The standard minimal has been introduced in 1969 by Lewis et
al.~\cite{LewiGoodMill69}.
It is a linear congruential generator driven by the recurrence
\[
  x_{n+1} = 16807x_n \mod (2^{31}- 1).
\]

\subsection{Multiple Recursive Generator}

These generators are based on the recurrence
\[
  {x_n} = (a_1 x_{n-1} + \cdots + a_k x_{n-k}) \mod {m}.
\]
The choice of the parameters is discussed by \citep{LEcu99}, who provides the code for the MRG32k3a generator, proposed in \oratio{}.

\subsection{Generators in $\mathcal{F}_2$}

\subsubsection{Mersenne twister generator}

The Mersenne twister is a pseudorandom number generator developed in
1997 by Makoto Matsumotoand Takuji Nishimura~\cite{MatsNish98}.
While it is reputated to be robust and fast, it is not exempt of defaults,
since if fails some tests defined in TestU01 (L'Ecuyer and
Simard~\cite{LEcuSima07}).
The state length is a vector of 624 integers.

\begin{verbatim}
\end{verbatim}

\section{Non-uniform random numbers generation}

While uniform random number generation are of deep importance, the
practician will mainly look for other distributions.
The main reference for non-uniform random number generation is the
book written by Luc Devroye~\cite{Devr86}, that can be
downloaded free of charge from the author's website, at the address
\url{http://cg.scs.carleton.ca/~luc/rnbookindex.html}.
A copy can also be found at the adress
\url{http://www.iro.umontreal.ca/~bastin/books/nonuniformrandomvariates.zip}.
We have also used Evans, Hastings and Peacock~\cite{EvanHastPeac93} as
reference.

The library aims to use a similar naming scheme amongst the various
distributions.
Each routine name will start by \verb|ran_|, followed by the
distribution name and the function of the routine.

\subsection{Univariate distributions}

\subsubsection{Discrete distributions}

%%%%%%%%%%%%%% BERNOUILLI DISTRIBUTION %%%%%%%%%%%%%%

\paragraph{Bernouilli distribution}
\index{distribution!Bernouilli}

A Bernouilli variable can take two value only: 1 (success) or 0
(failure). The probability of success, denoted by $p$, is the
Bernouilli parameter; $0 \leq p \leq 1$.

You can draw some value from a Bernouilli distribution with parameter $p$
using the function:
\begin{verbatim}
int ran_bernouilli(const Random *random, double p);
\end{verbatim}
You can also create a generator more specifically designed for the
Bernouilli distribution with the following function:
\begin{verbatim}
Random *ran_bernouilli_new(unsigned long *seed, int n, double p);
\end{verbatim}
A realization can then be obtained by using the function
\begin{verbatim}
int ran_bernouilli_get_val(const Random *random);
\end{verbatim}
The parameter $p$ can be retrieved with
\begin{verbatim}
double ran_bernouilli_get_param(const Random *random);
\end{verbatim}
On the other side, you can define the parameter using
\begin{verbatim}
double ran_bernouilli_set_param(Random *random, double p)
\end{verbatim}

%%%%%%%%%%%%%%%% GEOMETRIC DISTRIBUTION %%%%%%%%%%%%%%%%

\paragraph{Geometric distribution}

$X$ is geometrically distributed with parameter $p \in (0,1)$ when
\[
P[X = i] = p(1-p)^{(i)}, \ i = 0,1,\ldots
\]
A for the Bernouilli variate, $p$ is the success probability of one
trial, so that $0 \leq p \leq 1$.
$i$ is number of failures until the first success.
The geometric distribution is important because it is the distribution
of the waiting time until success in a sequence of Bernouilli
trials. The geometric distribution is also sometimes called the Pascal
distribution.

You can obtain a draw from a geometric distribution of parameter
$p$ with a general Random object by using the function
\begin{verbatim}
int ran_geometric(const Random *random, double p)
\end{verbatim}
The geometric distribtion object can be created with the function
\begin{verbatim}
Random *ran_geometric_new(unsigned long *seed, int n, double p)
\end{verbatim}
A random number can then be drawn with the function
\begin{verbatim}
int ran_geometric_get_val(const Random *random)
\end{verbatim}
The parameter $p$ can be retrieved and set with the functions
\begin{verbatim}
double ran_geometric_get_param(const Random *random)
\end{verbatim}
\begin{verbatim}
void ran_geometric_set_param(Random *random, double p)
\end{verbatim}

%%%%%%%%%%%%%% RANDOM INTEGERS GENERATION %%%%%%%%%%%%%%

\paragraph{Integer numbers}

A random integer variable $X$, associated to the parameter $max$, is
assumed to be a discrete variable that presents the property that
\[
P[X = i] = \frac{1}{\max},\ \forall i \in [0, \max[ \ \cap\ \mathcal{N}.
\]
\begin{center}
Variate $X$: $\max$.
Range $\max \in \mathcal{N}_0.$
\end{center}

A draw can be generated with the function
\begin{verbatim}
void ran_intrandom_set_max(Random *ird, unsigned long max)
\end{verbatim}

An integer random generator can be created with the function
\begin{verbatim}
Random *ran_intrandom_new(unsigned long *seed, int n,
                          unsigned long max)
\end{verbatim}

\begin{verbatim}
unsigned long ran_intrandom_get_val(const Random *ird)
\end{verbatim}

The maximum value can be retrieved with the function
\begin{verbatim}
unsigned long ran_intrandom_get_max(const Random *ird)
\end{verbatim}

\subsubsection{Continuous distributions}

%%%%%%%%%%%%%% ARCSINE DISTRIBUTION %%%%%%%%%%%%%%

\paragraph{Arcsine distribution}

The arcsine distribution is a special case of the beta distribution with both parameters equal to $1/2$.
The density function is:
\[
f(x) =
\begin{cases}
\frac{1}{\pi\sqrt{x(1-x)}} & \text{if } 0\leq x \leq 1,\\
0 & \text{ otherwise}.
\end{cases}
\]

You can pick a number from the arcsine with the function
\begin{verbatim}
double ran_arcsine(const Random *random)
\end{verbatim}

It is also possible to create an object more specifically designed for the arcsine distribution with the function
\begin{verbatim}
Random *ran_arcsine_new(unsigned long *seed, int n)
\end{verbatim}
and then use this function to get some value:
\begin{verbatim}
double ran_arcsine_get_val(const Random *random)
\end{verbatim}

%%%%%%%%%%%%%% CAUCHY DISTRIBUTION %%%%%%%%%%%%%%

\paragraph{Cauchy distribution}

\begin{verbatim}
double ran_cauchy(const Random *random,
                  double a, double b)
\end{verbatim}
\begin{verbatim}
Random *ran_cauchy_new(unsigned long *seed, int n
                       double a, double b)
\end{verbatim}
\begin{verbatim}
double ran_cauchy_get_val(const Random *random)
\end{verbatim}
\begin{verbatim}
double ran_cauchy_standard_get_val(const Random *random)
\end{verbatim}
\begin{verbatim}
double ran_cauchy_get_parameters(const Random *random,
                                 double *a, double *b)
\end{verbatim}
\begin{verbatim}
double ran_cauchy_set_parameters(Random *random,
                                 double a, double b)
\end{verbatim}

%%%%%%%%%%%%%% ERLANG DISTRIBUTION %%%%%%%%%%%%%%

%\paragraph{Erlang distribution}
%\index{distribution!erlang}

%%%%%%%%%%%%%% EXPONENTIAL DISTRIBUTION %%%%%%%%%%%%%%

\paragraph{Exponential distribution}
\index{distribution!exponential}

The exponential distrbution, described in
Section~\ref{sec:stat_exponential}, can be generated with the
function
\begin{verbatim}
double ran_exponential(const Random *random, double lambda)
\end{verbatim}
where $\lambda$ is the distribution parameter.
The other functions are
\begin{verbatim}
Random *ran_exponential_new(unsigned long *seed, int n, double lambda)
double ran_exponential_get_val(const Random *random)
double ran_exponential_get_parameters(const Random *random)
void ran_exponential_set_parameters(Random *random, double lambda)
\end{verbatim}

%%%%%%%%%%%%%% GAMMA DISTRIBUTION %%%%%%%%%%%%%%

\paragraph{Gamma distribution}

\[
f(x) =
\begin{cases}
\frac{1}{\Gamma(a) b^a} x^{a-1} e^{-\frac{x}{b}} & \text{if } x > 0, \\
0 & \text{otherwise}.
\end{cases}
\]

\begin{verbatim}
double ran_gamma(const Random *random, double a, double b)
\end{verbatim}
\begin{verbatim}
Random *ran_gamma_new(unsigned long *seed, int n,
                      double a, double b)
\end{verbatim}
\begin{verbatim}
double ran_gamma_get_val(const Random *random)
\end{verbatim}
\begin{verbatim}
double ran_gamma_get_parameters(const Random *random,
                                double *a, double *b)
\end{verbatim}
\begin{verbatim}
double ran_gamma_set_parameters(Random *random,
                                double a, double b)
\end{verbatim}

%%%%%%%%%%%%%% GUMBEL DISTRIBUTION %%%%%%%%%%%%%%

\paragraph{Gumbel distribution}

You can obtain a value from a Gumbel distribution with
parameters $a$, $b$, using the function:
\begin{verbatim}
double ran_gumbel(const Random *random, double a, double b)
\end{verbatim}

You can create a generator more specifically designed for the
Gumbel distribution with the following function:
\begin{verbatim}
Random *ran_gumbel_new(unsigned long *seed, int n,
                       double a, double b)
\end{verbatim}
Then a value can be obtained by using the function:
\begin{verbatim}
int ran_gumbel_get_val(const Random *random)
\end{verbatim}
The parameters $a$, $b$ can be retrieved with:
\begin{verbatim}
void ran_gumbel_get_parameters(const Random *random,
                               double *a, double *b)
\end{verbatim}
On the other side, you can define the parameter using:
\begin{verbatim}
void ran_gumbel_set_parameters(Random *random,
                               double a, double b)
\end{verbatim}

%%%%%%%%%%%%%% JOHNSON DISTRIBUTION %%%%%%%%%%%%%%

\paragraph{Johnson distribution}

The Johnson families are part of the most flexible ones, and are able
to mimic most of standard distributions.
Its general form is 
\[
f(x) = \frac{\delta}{\lambda\sqrt{2\pi}} g' \left(
  \frac{x-\xi}{\lambda} \right)
e^{ -\frac{\left(\gamma + \delta g \left( \frac{x-\xi}{\lambda} \right)\right)^2}{2} },
\]
where $g$ is one of the following four transformations:
\[
g(y) = \begin{cases}
\ln y & \mbox{lognormal family} \\
\sinh^{-1}(y) = \ln \left( y + \sqrt{y^2+1}\right) & \mbox{unbounded family} \\
\ln \left( \frac{y}{1-y} \right) & \mbox{bounded family} \\
y & \mbox{normal family}.
\end{cases}
\]
\begin{center}
Variate $J$: $\xi$, $\lambda$, $\gamma$, $\delta$.\\
Localisation parameter $\xi$.\\
Scale parameter $\lambda$.
Form parameters: $\gamma$, $\delta$. \\
\end{center}
The support of the distribution depends on the transformation $g$. It
is $(-\infty,\,\infty)$ for the unbounded and normal family,
$[\xi,\,\xi + \lambda]$ for the bounded family, and $[\xi, \infty)$
for the lognormal family.

The associated cumulative distribution function is:
\[
F(x) = \Phi\left(\gamma+\delta g \left( \frac{x-\xi}{\lambda} \right)\right).
\]
Since the normal and lognormal families reduce to normal and lognormal
random variables, we consider only bounded and unbounded families.
The bounded family will be identified with the letters \verb|SB|,
while the unbounded family will be identified with the letters \verb|SU|.

\subparagraph{Unbounded family}

A draw from an unbounded Johnson variate of parameters $\xi$,
$\lambda$, $\gamma$, $\delta$ can be obtained by calling the function
\begin{verbatim}
double ran_johnsonSU(const Random *random,
                     double xi, double lambda,
                     double gamma, double delta);
\end{verbatim}

As usual, you can define a specific generator using the function
\begin{verbatim}
Random *ran_johnsonSU_new(unsigned long *seed, int n,
                          double xi, double lambda,
                          double gamma, double delta);
\end{verbatim}
and then draw on it with repeated call to
\begin{verbatim}
double ran_johnsonSU_get_val(const Random *random);
\end{verbatim}

The parameters can be retrieved with the function
\begin{verbatim}
void ran_johnsonSU_get_parameters(const Random *random,
                                  double *xi, double *lambda,
                                  double *gamma, double *delta);
\end{verbatim}
and set with
\begin{verbatim}
void ran_johnsonSU_set_parameters(Random *random,
                                  double xi, double lambda, 
                                  double gamma, double delta);
\end{verbatim}

\subparagraph{Bounded family}

The functions are the same as for the unbounded case, except that the
letters \verb|SU| in function names have to be replaced by \verb|SB|.
This results in the following set of functions:
\begin{verbatim}
double ran_johnsonSB(const Random *random,
                     double xi, double lambda,
                     double gamma, double delta);
\end{verbatim}
\begin{verbatim}
Random *ran_johnsonSB_new(unsigned long *state,
                          double xi, double lambda,
                          double gamma, double delta);
\end{verbatim}
\begin{verbatim}
double ran_johnsonSB_get_val(const Random *random);
\end{verbatim}
\begin{verbatim}
void ran_johnsonSB_get_parameters(const Random *random,
                                  double *xi, double *lambda,
                                  double *gamma, double *delta);
\end{verbatim}
\begin{verbatim}
void ran_johnsonSB_set_parameters(Random *random,
                                  double xi, double lambda, 
                                  double gamma, double delta);
\end{verbatim}

%%%%%%%%%%%%%% LOGNORMAL DISTRIBUTION %%%%%%%%%%%%%%

\paragraph{Lognormal distribution}

The density function of a log-normal variate $X$ is
\[
f(x) =
\begin{cases}
\frac{1}{x \sqrt{2\pi} \sigma} e^{ -\frac{ ( \log x - \mu )^2
}{2\sigma^2} } & \text{ if } x > 0, \\
0 & \text{ otherwise.}
\end{cases}
\]
We then write $X \sim LN(\mu, \sigma^2)$. The distribution function is
\[
F(x) =
\begin{cases}
\frac{1}{2} \left( 1 + \text{erf} \left( \frac{\ln x -
 \mu}{\sqrt{2}\sigma} \right) \right)
 & \text{ if } x > 0, \\
0 & \text{ otherwise.}
\end{cases}
\]

$X \sim LN(\mu, \sigma^2)$ if and only if $\ln X \sim N(\mu, \sigma^2)$.

\begin{verbatim}
double ran_lognormal(const Random *random,
                     double mu, double sigma)
\end{verbatim}

\begin{verbatim}
double ran_lognormal_icdf(const Random *random,
                          double mu, double sigma)
\end{verbatim}

\begin{verbatim}
Random *ran_lognormal_new(unsigned long *seed, int n,
                          double mu, double sigma)
\end{verbatim}

\begin{verbatim}
double ran_lognormal_get_val(Random *random)
\end{verbatim}

\begin{verbatim}
double ran_lognormal_get_val_with_parameters
  (Random *random, double mu, double sigma)
\end{verbatim}

\begin{verbatim}
void ran_lognormal_get_parameters
  (const Random *random, double *mu, double *sigma)
\end{verbatim}

\begin{verbatim}
void ran_lognormal_set_parameters(Random *random,
                                  double mu, double sigma)
\end{verbatim}

%%%%%%%%%%%%%%% NORMAL DISTRIBUTION %%%%%%%%%%%%%%%

\paragraph{Normal distribution}

A random variable follows a normal distribution $N(\mu, \sigma^2)$ if its density function is given by:
\[
f(x) = \frac{1}{\sqrt{2\pi}\sigma}e^{-\frac{(x-\mu)^2}{2\sigma^2}}.
\]

It is possible to use a generic object of type \verb|random| to get
values from a normal distribution with the following function:
\begin{verbatim}
double ran_normal(const Random *random, double mu, double sigma)
\end{verbatim}
The method is based on the CDF inversion, which can also be called
using the function
\begin{verbatim}
double ran_normal_icdf(const Random *random, double mu, double sigma)
\end{verbatim}

It can however be preferable to define a random generator dedicaced to
deliver random values, with the function
\begin{verbatim}
Random *ran_normal_new(unsigned long *seed, int n,
                       double mean, double stdv)
\end{verbatim}
A value can then be obtained with the function
\begin{verbatim}
double ran_normal_get_val(Random *random)
\end{verbatim}
Is it also possible to get values from an normal $N(\mu,\sigma^2)$
with the function
\begin{verbatim}
double ran_normal_get_val_with_parameters(Random *random,
				          double mean, double sigma)
\end{verbatim}
where $\mu$ is the variance and $\sigma$ is the standard deviation.

Various alternatives has been proposed to generate normal numbers.
The most famous one has been proposed by Box and Muller, and transform two uniform numbers on (0,1) to $N(0,1)$ numbers.
It is mathematically elegant, but in practice, slower than the inversion approach, and not compatible with variance-reduction technique.
A second alternative is the acceptance-rejection method from the Laplace distribution, but it is even slower.
These two different generation methods are present in \oratio{} for illustration purpose only, but should but used for serious purposes.
The file \verb|normal_generation.c|, in the examples directory, compare these methods for the generation of $10^8$ draws from a Gaussian distribution.
On a AMD FX(tm)-8120 Eight-Core Processor, we obtained the following results:
\begin{verbatim}
Time with acceptance-rejection method, for 100000000 numbers: 27.978 seconds
Time with Box-Muller method, for 100000000 numbers: 14.276 seconds
Time with inversion method, for 100000000 numbers: 9.121 seconds
\end{verbatim}

\paragraph{Student distribution}

\begin{description}
\item[Density function]
\[
f(x) =
\begin{cases}
\frac{1}{2} + & \mbox{if $\nu$ is odd} \\
\frac{1}{2} + & \mbox{if $\nu$ is even} 
\end{cases}
\]
\end{description}

\paragraph{Triangular distribution}

A triangular distribution generator can be obtained with the function
\begin{verbatim}
Random *ran_triangular_new(unsigned long *seed, int n,
                           double x_min, double x_max, double mode)
\end{verbatim}

\begin{verbatim}
double ran_triangular(const Random *random,
		      double x_min, double x_max, double mode)
\end{verbatim}

\begin{verbatim}
double ran_triangular_get_val(const Random *random)
\end{verbatim}

\begin{verbatim}
void ran_triangular_get_parameters(const Random *random,
				   double *x_min, double *x_max, double *mode)
\end{verbatim}

\begin{verbatim}
void ran_triangular_set_parameters(Random *random,
				   double x_min, double x_max, double mode)
\end{verbatim}

\paragraph{Truncated normal distribution}

The truncated normal distribution is the probability distribution of a
normally distributed random variable whose value is either bounded
below or above (or both).
Assume that $Y ~ N(\mu, sigma^2)$.
\begin{description}
\item[Parameters]
\begin{itemize}
\item $\mu$: mean of $Y$;
\item $\sigma$: standard deviation of $Y$;
\item $a$, $b$: lower and upper bounds of $X$, respectively.
\end{itemize}
\item[Density function]
\[
f(x) = \frac{\frac{1}{\sigma} \phi \left( \frac{x-\mu}{\sigma} \right)}
{\frac{1}{\sigma} \Phi \left( \frac{b-\mu}{\sigma} \right) - \Phi \frac{1}{\sigma}\left( \frac{a-\mu}{\sigma} \right)}
\]
\end{description}

\paragraph{Truncated normal distribution with mass points}

\subsection{Multivariate distributions}

Function dealing with multivariate distributions are similare to the
univariate case, except the presence of the prefix \verb|m_| in from
of the distribution name.

\subsubsection{Multivariate normal}

The multivariate normal is described in \ref{sec:st_mvn}, and can be generated with the following functions.
A draw from a standard multivariate normal of dimension $n$ can be obtained with the function
\begin{verbatim}
void ran_mvn_gaussian(Random *ran, int n, double *x, int incx)
\end{verbatim}
As usual, \verb|ran| is a $U(0,1)$ random numbers generator, and the draw in stored in the $n$-dimensional vector $x$, using an increment of \verb|incx|.

A Box-Muller variant is given by the function
\begin{verbatim}
void ran_mvn_gaussian_box_muller(Random *ran, int n,
                                 double *x, int incx)
\end{verbatim}
It is advised to not use this variant except for testing purposes, as it is slower than the previous function, based on inversion of the distribution function. Moreover, the Box-Muller approach is not compatible with standard variance reduction techniques.

A draw from a multivariate normal $\cN(\bmu, \bSigma)$, where $\bmu$ is the mean vector and $\bSigma$ is the variance-covariance matrix, can be obtained from the function
\begin{verbatim}
void ran_mvn(Random *ran,
             const enum CBLAS_ORDER Order, int n, double *L, int ldL,
             double *mu, int incmu, double *x, int incx)
\end{verbatim}
In this function $L$ is a lower triangular matrix of leading dimension \verb|ldL|, with orientation given by \verb|Order|, such that $LL' = \bSigma$.


\subsubsection{Unit sphere}

Let denote by $C^m$ the n-dimensional unit sphere, that is
\[
C^n = \lbrace x \in \mathcal{R}^m \mbox{ s.t. } \|x\|_2 = 1 \rbrace.
\]
The following function allows the generation of a random vector
uniformly distributed on $C^n$.

A specific uniform random number generator on $C^m$ can be created
with the function
\begin{verbatim}
Random *ran_m_unit_sphere_new(unsigned long *seed, int n, int m)
\end{verbatim}
which than allow to draw uniformly on $C^m$ using the function
\begin{verbatim}
void ran_m_unit_sphere_get_val(Random *ran, int m, double *x)
\end{verbatim}
that stores the drawn vector in $x$.
If you prefer to use a generic random numbers generator, the following
function achieves the same effect:
\begin{verbatim}
void ran_m_unit_sphere(Random *ran, int m, double *x)
\end{verbatim}

%\subsection{Generic techniques}

\subsection{Acceptance-rejection}

\begin{algo}{Acceptance-rejection technique}

\mbox{}

\noindent
Repeat
\begin{enumerate}
\item
Generate $Y$ with density $r(x)$.
\item
Generate $U: U (0, 1)$ independent of $Y$.
\end{enumerate}
until $U \leq f (Y )/t(Y )$; $\qquad$ return Y.
\end{algo}

\subsection{Special cases}

\subsubsection{Ratio of distributions}
\index{distribution@ratio}

It is sometimes useful to compute the ratio of two random
distributions, which can be performed with the following functions.
\begin{verbatim}
void ran_generate_ratio_distribution(Random *ran1, Random *ran2,
			             int n, double *x1, int incx1,
				     double *x2, int incx2,
				     double *r, int incr)
\end{verbatim}
Both random numbers generators \verb|ran1| and \verb|ran2| must of
course be initialized, and producing independent sequences (except in
very specific cases). \verb|n| designs the number of draws to
produces.
The vectors $x_1$ will contain draws from the first distribution,
using \verb|ran_1|, using the increment \verb|incx1| (see
Section~\ref{sub:laconventions}), while $x_2$ will contain draws
produced by \verb|ran_2|, using the increment \verb|incx2|.
The ratios between these draws will be stocked in the vector $r$,
using the increment \verb|incr|.

This function can of course be memory consuming as it stores vectors
\verb|x1|,  \verb|x2| and  \verb|r|. A simpler function is
\begin{verbatim}
inline double ran_generate_ratio(Random *ran1, Random *ran2,
				 double *x1, double *x2)
\end{verbatim}
This function simply draws one realization from \verb|ran1|, one from
\verb|ran2|, stores them in elements \verb|x1| and \verb|x2|
respectively, and returns their ratio.

\subsubsection{Truncated distributions}

For any random variable $X$, it is possible to truncate its
distribution produced to some interval $[a,b]$ when inversion is used
as the generation technique.
This can be obtained by setting three parameters of the random
generator object:
\begin{itemize}
\item
the desired lower bound of the distribution function, i.e. $F(a)$;
\item
the desired upper bound of the distribution function, i.e. $F(b)$;
\item
the probability mass of $a$, i.e. $P[X=a]$.
\end{itemize}
These parameters can be set with the function
\begin{verbatim}
void ran_random_set_truncation(Random *random, double l,
                               double h, double p);
\end{verbatim}
$l$ and $h$ correspond to lower and upper bounds respectively, while
$p$ is the value $P[X=a]$.
By default, $l = 0$, $h = 1$, and $p = 0$, so that no truncation is performed.

In order to illustrate how the truncation is performed, let
$\tilde{F}$ denote the distribution $F$ truncated to the interval
$(a,b]$, i.e.,
\[
 \tilde{F}(x) = \frac{F(x)-F(a)}{F(b)-F(a)}
\]
for $a \le x\le b$, $\tilde F(x) = 0$ for $x \le a$, and $F(x) = 1$
for $x\ge b$.
Let $U \sim U(F(a),\, F(b))$ and $X=F^{-1}(U)$.
Since $F$ is monotone, we have
\begin{multline*}
P[X\leq x] = P[U\leq F(x)] =\\
\begin{cases}
  0 & \mbox{if } F(x)\le F(a) (\mbox{or if } x \le a), \\
  \frac{F(x)-F(a)}{F(b)-F(a)} & \mbox{if } F(a) \le F(x) \leq F(b)
  (\mbox{or if } a \le x\leq b), \\
  1 & \mbox{if } F(x) \ge F(b)  (\mbox{or if } x \ge b).
\end{cases}
\end{multline*}
Therefore, $X$ has distribution $\tilde{F}$.

If $F$ has a jump of size $p$ at $a$ and we want to truncate to
$[a,b]$, then the algorithm should return $a$ with probability $p$.
To achieve that, it suffices to generate $U$ uniform over $(F(a)-p,\,
F(b))$.

\section{Stochastic processes}

\begin{defn}{Stochastic process}
Given a probability space $(\Omega, \cF, P)$ and a measurable space
$(S, \Sigma)$ an $S$-valued \Index{stochastic process} $X$ is a collection of $S$-valued random variables on $\Omega$ indexed by a totally ordered set $T$:
\[
X = \lbrace X_t\,|\, t \in T\rbrace,
\]
where each $X_t$ is $S$-valued random variable on $\Omega$.
The space $S$ is then called the state space of the process.
\end{defn}

\subsection{One-dimensional Brownian motion}

\begin{defn}{Standard one-dimensional Brownian motion}
\index{Brownian motion}
A standard one-dimensional Brownian motion on $[0,T]$ is a stochastic
process $\lbrace W(t),\ 0 \leq t \leq T \rbrace$ with the following properties
\begin{enumerate}[(a)]
\item
 W(0) = 0;
\item
 the mapping $t \rightarrow W(t)$ is a continuous function on $[0,T]$, almost surely;
\item
 the increments
\[
 \lbrace W(t_1) - W(t_0), W(t_2) - W(t_1),\ldots, W(t_k) - W(t_{k-1}), \rbrace
\]
are independent for any $k$ and any $0 \leq t_0 < t_1 < \ldots < t_k \leq T$;
\item
 $W(t)-W(s) \sim N(0, t-s)$, for any $0 \leq s < t \leq T$.
\end{enumerate}
\end{defn}

For constants $\mu$ and $\sigma > 0$, a process $X(t)$ is a Brownian motion
with drift $\mu$ and diffusion coefficient $\sigma^2$ (abbreviated
$X \sim BM(\mu, \sigma^2)$) is
\[
\frac{X(t)-\mu t}{\sigma}
\]
is a standard Brownian motion $W$.
Therefore, we may construct $X$ from a standard Brownian motion $W$ by setting
\[
 X(t) = \mu t + \sigma W(t).
\]
It follows $X(t) \sim N(\mu t, \sigma^2 t)$, and $X$ solves the stochastic
differential equation 
\[
 dX(t) = \mu dt + \sigma dW(t).
\]
We can also start from a non-zero value $X(0) = \alpha$ by simply adding $x$
to each $X(t)$.

\subsection{One-dimensional geometric Brownian motion}

$S$ follows a geometric Brownian motion of drift $\mu$ and volatility parameter $\sigma$ if it satisfies the stochastic differential equation
\[
dS_t = \mu S_t dt + \sigma S_t dW_t,
\]
or, equivalently,
\[
S(t) = S(0)e^{(\mu-\frac{1}{2}\sigma^2)t+\sigma W(t)}.
\]

\subsection{Multiple-dimensional geometric Brownian motion}

Suppose that we have a correlation matrix $C$:
\[
C =
\begin{pmatrix}
1 & \rho_{1,2} & \ldots & \rho{1,n} \\
\rho_{2,1} & 1 & \ldots & \rho{2,n} \\
\vdots & \vdots & \ddots & \vdots \\
\rho_{n,1} & \rho_{n,2} & \ldots & 1
\end{pmatrix}
\]
$\rho_{ij}$ is the correlation between $X_i(t)$ and $X_j(t)$.
The corresponding covariance matrix is denoted $\Sigma$:
\[
\Sigma_{ij} = \sigma_i \sigma_j \rho_{ij}.
\]
Assume that $A$ is any matrix for which
\[
AA^T = C.
\]
For instance, $A$ can be chosen as the Cholesky factor of $C$.

\section{Random permutations}

A random permutation\index{permutation!random}, or
shuffling\index{shuffling|see{permutation}}, is a random ordering of a
set of objects, that is, a permutation-valued random variable.
There are two basic algorithms for doing this, both popularized by
Donald Knuth.
The first is simply to assign a random number to each object, and then
to sort the cards in order of their random numbers. This will generate
a random permutation, unless two of the random numbers generated are
the same. This can be eliminated either by retrying these cases, or
reduced to an arbitrarily low probability by choosing a sufficiently
wide range of random number choices.

The second, generally known as the Knuth shuffle or Fisher-Yates
shuffle, is a linear-time algorithm which involves moving through the
objects set, swapping each object in turn with another object from a
random position in the part of the set that has not yet been passed
through (including itself).
This is the implemented method, that comes in two versions, one for
generic pointers arrays, on for double vectors:
\begin{verbatim}
void ran_shuffle_generic(Random *r, int n, void *x, int incx);
void ran_shuffle_double(Random *r, int n, double *x, int incx);
\end{verbatim}
In these functions, \verb|r| is an initialized random numbers
generator, $n$ is the number of elements to shuffle, stored or pointed
by the vector $x$, which has an increment \verb|incx|.

\section{Input-ouput functions}

Two simple functions are available to write random draws in an output
file $\texttt{f}$.
The first function allows to write a random vector $v$ of size $d$.
\begin{verbatim}
void ran_write_vector(FILE *f, double *v, int d);
\end{verbatim}
This function is generalized to an array of dimensions $m \times d$,
and leading dimension $\texttt{ldS}$ (see
Section~\ref{sub:laconventions}).
\begin{verbatim}
void ran_export_sample(FILE *f, int m, int d, double *S, int ldS);
\end{verbatim}
